\chapter{Conclusion and open problems}\label{ch:conclusion}

This thesis asked whether safety in feedback systems with adaptive agents requires knowing who the agents are. For the systems studied here, it does not. Prices reveal the agents only through a merged measure of their memories and weights, and recovering the agents from that measure costs trading capital of order $\delta^{-4}$, experiments of order $(\Delta\eta)^{-(4K-2)}$, and with general memory an accuracy that improves only logarithmically (Chapter~\ref{ch:p1}). Safety depends on the agents only through their total feedback weight, which is certified at the parametric rate under a memory bound and cannot be certified without one; a guard built on that certificate is the largest safe admission rule (Chapter~\ref{ch:p2}). The framework of Chapter~\ref{ch:framework} places these results in one picture, in which the cost of a question is set by its behaviour at long memory. Chapter~\ref{ch:beyondlinear} shows that the safety bound needs only the agents' total sensitivity, even when their rules are nonlinear and their memories and weights change, and that learning makes certificates perishable at a computable rate. In data, disclosed trader classes reveal what prices do not, and the safety statistic forecasts adverse excursions in thin books, while the mechanism behind that forecast remains unidentified.

The policy message is direct. Supervisors who want to control transient risk should ask for the information that certifies safety, a bound on memory and on total sensitivity, renewed as often as the agents learn, rather than for the identities of the agents. Identity-based regimes may serve other purposes, such as enforcement and attribution, but this thesis shows they are neither necessary nor efficient for safety.

\section*{Open problems}

\begin{enumerate}
\item \emph{The matching lower bound.} Prove Theorem~\ref{thm:dich}(iv), so that the cost of every regular linear question is known up to constants, and extend the dichotomy to non-linear questions such as the peak response.
\item \emph{Minimal disclosure for safety.} Characterize the least informative disclosure under which a given safety question becomes certifiable from finite data. Theorem~\ref{thm:dich}(iii) says a memory bound is necessary for the summed response; is it sufficient in a sense that can be made operational?
\item \emph{Strategic agents.} Prove or refute Conjecture~\ref{conj:strategic}.
\item \emph{Observational certification.} Develop finite-sample certificates from observational order-book data with dependence and non-stationarity, replacing the bootstrap diagnostics of Chapter~\ref{ch:p2}.
\item \emph{Sharp multi-asset and multi-state bounds.} Replace the induced-norm bound of Proposition~\ref{prop:multiasset} with sharp bounds for positive cross-impact and for grid models with governor dynamics.
\item \emph{The mechanism of thin-book risk.} Run Tests E1 and E2 of Chapter~\ref{ch:evidence}.
\end{enumerate}
